\section{Discussion}
\label{sec:s2-discussion}

\subsection{The useful unit is a mediated action, not an ACK}

The focal replay changes the interpretation of ACK assistance. A standalone
ACK is not a receiver update; it is an action on sender information state. Its
receiver-side effect is mediated through subsequent DATA and contention. This
is why positive confirmation lead coexists with adverse true AoI at N5 and
adverse formation loss at N10. A decision rule based only on expected ACK
delay, ACK success probability, or age reduction is structurally incomplete.
It must also account for how the changed belief alters future update demand.

The exact airtime identity makes the same point at policy level. Standalone
feedback reduces offered load only when DATA airtime saved exceeds additional
ACK airtime. Under N5 CSMA, the earlier EXP14E audit recovers only 49.1\% of
direct ACK cost through DATA savings; at N10 it induces rather than saves DATA.
Under ALOHA, DATA savings cover 247.4\% of direct ACK airtime. These signs are
reproduced on new EXP14I traces and survive the formal within-seed interaction
test in EXP16.

\subsection{Why the MAC reverses feedback value}

Carrier-sensing service and collision-only ALOHA turn the same confirmation
into different future opportunity costs. Under CSMA, piggyback can transport
cumulative feedback on already useful DATA while standalone frames consume
service opportunities and add collision exposure. At N5/N10, that direct cost
is not repaid. Under the tested ALOHA load, delayed confirmation sustains more
redundant DATA and collisions; a fast standalone path reduces both enough to
repay its own traffic.

Equation~\eqref{eq:s2-mac-selector} is intentionally categorical. The data do
not support a finely fitted state predictor: the focal study contains only 20
independent decisions in each supported cell and heterogeneous signs. The
selector instead encodes a mechanism visible before runtime and then receives
an entirely disjoint directional validation followed by a matched factorial
interaction test. Its exact alias gate also prevents a semantic
sleight of hand: within a cell, the candidate is just the chosen fixed route.

The EXP16 interaction is the strongest evidence for the MAC-conditioned
mechanism, but its interpretation remains deliberately narrow. The configured
CSMA/ALOHA factor changes abstract access semantics while holding the random
loss realization fixed; it does not identify which detailed 802.11 feature
would reproduce the effect. Moreover, adaptive and fixed-deadline hybrid are
statistically unresolved under ALOHA. The evidence therefore favors retaining
a standalone feedback path in that cell, not the universal optimality of the
specific adaptive scheduler.

\subsection{Why configured MAC type is not a sufficient selector}

EXP17 distinguishes a mechanism from an operating policy. The base-cell
interaction is reproducible, but background occupancy, burst severity, swarm
size, and reverse-path asymmetry change both the opportunity cost of ACKs and
the DATA actions that confirmation suppresses. In loaded cells the route
effect can vanish because both routes approach the same unsafe service limit.
At N10 the CSMA resource benefit may persist while the ALOHA control benefit
does not. Reverse asymmetry can improve RMSE while increasing offered airtime.

Consequently, the categorical rule in \eqref{eq:s2-mac-selector} is a compact
description of the supported base regime, not a deployable selector. A future
policy would require observable load/channel/size state and an independently
preregistered validation. Reusing EXP17 to tune and then claim confirmation on
the same contexts would erase the evidence boundary, so no post-hoc repair is
performed here.

\subsection{Support failure is a result}

At N20 the selected adaptive policy already behaves almost piggyback-only.
Only 11 of 300 development trajectories contain an eligible post-transient
standalone decision. Repeatedly adding seeds or moving the analysis window
would change the design after observing the support problem. Reporting the
failure preserves the estimand: EXP14G identifies event-conditional effects
only where the natural candidate process supplies decisions. Full-policy
comparisons remain possible, but they do not manufacture a conditional ACK
effect in unsupported states.

\subsection{Communication and control claims remain distinct}

Access scaling repairs catastrophic N20 contention, while the fixed load guard
shows that a lower channel demand can still worsen control and safety. The
formation result cannot therefore be inferred from a MAC metric alone.
Conversely, lower RMSE does not prove communication efficiency: the initial
Stressed default improves error while using more offered airtime than P10.
The paper reports both axes and retains raw safety counts rather than collapse
them into an unregistered scalar objective.

\subsection{Relation to the theoretical results}

Causal conservatism is an invariant: delayed or lost honest ACKs can make the
sender pessimistic but not optimistically fresher than receiver truth. The age
tail and formation ISS results are conditional on service minorization,
bounded motion, and an unsaturated linear region. They explain how a credible
service bound could propagate to safety, but the current generic Lyapunov bound
is too conservative for a useful numerical probability. Simulation verifies
protocol and mechanism contracts; it is not substituted for those theorem
assumptions.

The ACK-value proposition is similarly a decision envelope, not an empirical
guarantee that ACK transmission is beneficial. EXP14G supplies pathwise focal
contrasts on supported events and falsifies lead-as-value. It does not fit the
benefit-rate envelope required for a universal state-conditional certificate.
